\SourceLeaf{445}%
\[
(\xi\xi'+\eta\eta'+\zeta\zeta')\frac{\partial t}{\partial\theta}=\frac{\partial x}{\partial s}.\frac{\partial x}{\partial\theta}+\frac{\partial y}{\partial s}.\frac{\partial y}{\partial\theta}+\frac{\partial z}{\partial s}.\frac{\partial z}{\partial\theta},
\]
a quantity we shall denote by $u$, which is therefore itself a function of $\theta$ and $s$.

Differentiating with respect to $s$, we find
\[
\begin{aligned}
\frac{\partial u}{\partial s}={}&\frac{\partial\partial x}{\partial s^{2}}.\frac{\partial x}{\partial\theta}+\frac{\partial\partial y}{\partial s^{2}}.\frac{\partial y}{\partial\theta}+\frac{\partial\partial z}{\partial s^{2}}.\frac{\partial z}{\partial\theta}\\
&+\tfrac12\frac{\partial\left\{\left(\frac{\partial x}{\partial s}\right)^{2}+\left(\frac{\partial y}{\partial s}\right)^{2}+\left(\frac{\partial z}{\partial s}\right)^{2}\right\}}{\partial\theta}\\
={}&\frac{\partial\partial x}{\partial s^{2}}.\frac{\partial x}{\partial\theta}+\frac{\partial\partial y}{\partial s^{2}}.\frac{\partial y}{\partial\theta}+\frac{\partial\partial z}{\partial s^{2}}.\frac{\partial z}{\partial\theta},
\end{aligned}
\]
because
\[
\left(\frac{\partial x}{\partial s}\right)^{2}+\left(\frac{\partial y}{\partial s}\right)^{2}+\left(\frac{\partial z}{\partial s}\right)^{2}=1,
\]
and its differential is therefore $=0$.

But since all points [corresponding] to a constant $\theta$ lie on a shortest line, if $L$ denotes the zenith of the point to which $s,\theta$ correspond and the coordinates of $L$ are denoted by $X,Y,Z$, we have
\[
\frac{\partial\partial x}{\partial s^{2}}=\frac{X}{p},\qquad\frac{\partial\partial y}{\partial s^{2}}=\frac{Y}{p},\qquad\frac{\partial\partial z}{\partial s^{2}}=\frac{Z}{p},
\]
where $p$ denotes the radius of curvature. We therefore have
\[
p.\frac{\partial u}{\partial s}=X\frac{\partial x}{\partial\theta}+Y\frac{\partial y}{\partial\theta}+Z\frac{\partial z}{\partial\theta}=\frac{\partial t}{\partial\theta}(X\xi'+Y\eta'+Z\zeta').
\]
But
\[
X\xi'+Y\eta'+Z\zeta'=\cos L\lambda'=0,
\]
because $\lambda'$ evidently lies on the great circle whose pole is $L$. We therefore have
\[
\frac{\partial u}{\partial s}=0
\]
or $u$ is independent of $s$ and hence is a function of $\theta$ alone. But for $s=0$ we evidently have $t=0$, $\frac{\partial t}{\partial\theta}=0$, and hence $u=0$. We conclude that in general $u=0$, or
\[
\cos\lambda\lambda'=0.
\]
From this follows the beautiful theorem:
